\documentclass[11pt]{article}
\usepackage[T1]{fontenc}
\usepackage[utf8]{inputenc}
\usepackage{lmodern,amsmath,amssymb,amsthm,mathtools}
\usepackage[a4paper,margin=27mm]{geometry}
\usepackage{microtype,booktabs,array,longtable,enumitem,xcolor}
\usepackage{hyperref,listings,fancyhdr}
\hypersetup{colorlinks=true,linkcolor=blue!45!black,urlcolor=blue!45!black,citecolor=blue!45!black,
pdftitle={Euler's rotation theorem and Rodrigues' formula in iPRA},
pdfauthor={Research note prepared for Vladimir Reshetnikov}}
\pagestyle{fancy}\fancyhf{}\fancyhead[L]{\small Euler and Rodrigues in iPRA}\fancyhead[R]{\small Constructive scope and proofs}\fancyfoot[C]{\thepage}
\setlength{\headheight}{14pt}
\setlength{\emergencystretch}{2em}
\setlist{topsep=4pt,itemsep=3pt}
\lstset{basicstyle=\ttfamily\small,breaklines=true,columns=fullflexible,keepspaces=true,frame=single,rulecolor=\color{black!20}}
\newtheorem{theorem}{Theorem}[section]
\newtheorem{lemma}[theorem]{Lemma}
\newtheorem{corollary}[theorem]{Corollary}
\newtheorem{proposition}[theorem]{Proposition}
\theoremstyle{definition}\newtheorem{definition}[theorem]{Definition}
\newtheorem{remark}[theorem]{Remark}
\newcommand{\ipra}{\mathsf{iPRA}}
\newcommand{\Q}{\mathbb Q}\newcommand{\N}{\mathbb N}\newcommand{\R}{\mathbb R}
\newcommand{\tr}{\operatorname{tr}}\newcommand{\adj}{\operatorname{adj}}
\newcommand{\Id}{I}\newcommand{\SO}{\operatorname{SO}}
\newcommand{\norm}[1]{\left\lVert#1\right\rVert}
\newcommand{\abs}[1]{\left|#1\right|}
\newcommand{\crossmat}[1]{[#1]_{\times}}
\newcommand{\eps}{\varepsilon}
\newcommand{\PR}{\mathsf{PR}}
\title{\textbf{Euler's rotation theorem and Rodrigues' formula\\in intuitionistic primitive recursive arithmetic}\\[8pt]
\large Exact constructions, finite precision, and the obstruction to selecting an axis}
\author{Research note prepared for Vladimir Reshetnikov}
\date{November 14, 2026}
\begin{document}
\maketitle
\begin{abstract}
Rodrigues' formula has a direct constructive proof using fixed-size polynomial
identities. Euler's theorem needs a more careful statement. We prove an
unconditional exact theorem for rational proper orthogonal matrices, and an
exact theorem for primitive-recursive Cauchy matrices supplied with a numerical
bound separating them from the identity. Both constructions give the complete
Rodrigues normal form, including half-turns, and use only bounded arithmetic,
square roots with explicit moduli, and divisions with positive lower bounds.
For arbitrary represented real rotations we prove an unconditional
finite-precision axis theorem and an unconditional exact unit-quaternion
representation. We also give a complete obstruction: a uniform exact unit-axis
theorem for all such real rotations implies LLPO; a specific primitive-recursive
input family yields a non-realizable first-order arithmetic consequence.
Thus the unrestricted classical exact-axis assertion cannot be honestly
claimed as a theorem of the requested constructive theory.
The proofs are human-readable; the accompanying Lean files are an uncompiled
prototype and a detailed plan, not a completed formalization. Executable
supplements check exact rational certificates and polynomial identities.
\end{abstract}
\tableofcontents
\newpage

\section{What can, and cannot, be proved}
\label{sec:scope}

We use the standard matrix formulation of an orientation-preserving rigid
motion fixing the origin:
\begin{equation}\label{eq:so}
 R^T R=I,\qquad \det R=1,\qquad R\in\R^{3\times3}.
\end{equation}
A fixed point other than the origin is handled by replacing $x$ with $x-p$.
The orientation assumption matters: an arbitrary distance-preserving map
fixing a point can include a reflection. The classical Euler assertion is that
there is a unit vector $u$ with $Ru=u$, and that the restriction to its
perpendicular plane is a planar rotation. The requested Wikipedia pages
describe this assertion and the associated Rodrigues formula
\cite{eulerwiki,rodrigueswiki}; no proof step below depends on those expositions.

The matrix formulation itself is constructive. If a distance-preserving
map $F$ fixes $p$, put $g_i=F(p+e_i)-p$. Expanding the distances between
$p,p+e_1,p+e_2,p+e_3$ shows $g_i\cdot g_j=I_{ij}$. For
$y=F(p+x)-p$, the distances to $p$ and to $p+e_i$ give
$y\cdot g_i=x_i$ by the identity
$2a\cdot b=|a|^2+|b|^2-|a-b|^2$. Consequently
$|y-\sum_i x_i g_i|^2=0$, so $y=\sum_i x_i g_i$.
Thus $F(p+x)=p+Rx$ with columns $g_i$. Orientation preservation is
exactly $\det R=1$. Only the elementary sum-of-squares inference is used.

The phrase ``in $\ipra$'' requires a representation. First-order arithmetic has
natural-number variables, not an unrestricted type of real numbers. Throughout,
a represented real is a rational approximation program with a specified error
modulus. This distinction changes the exact existence theorem.

\begin{theorem}[Positive results proved in this note]\label{thm:results}
The following statements admit the primitive-recursive constructions and
intuitionistic arithmetic proofs described below.
\begin{enumerate}[label=(\roman*)]
\item \textbf{Rodrigues, without an extra hypothesis.} Given a unit vector
$u$ and real numbers $c,s$ satisfying $c^2+s^2=1$, the map
\begin{equation}\label{eq:rodmain}
 \mathcal R(u,c,s)x=cx+s(u\times x)+(1-c)(u\cdot x)u
\end{equation}
is a proper orthogonal map fixing $u$. On $u^\perp$ it is the planar rotation
specified by $(c,s)$. Taking $c=\cos\theta$, $s=\sin\theta$ gives the usual
formula, with primitive-recursive names for the trigonometric functions.
\item \textbf{Exact Euler for rational input.} From any rational matrix
satisfying \eqref{eq:so}, a primitive-recursive procedure constructs a nonzero
rational fixed vector $v$. Normalizing it gives a represented unit axis and
exact Rodrigues parameters. No nonidentity promise is required.
\item \textbf{Exact Euler for separated real input.} Put
$d=(3-\tr R)/2$. If an integer $h$ is supplied with $d\ge2^{-h}$, an explicit
primitive-recursive procedure constructs $u,c,s$ with
$u\cdot u=1$, $c^2+s^2=1$, and $R=\mathcal R(u,c,s)$. The fixed space is
exactly $\{tu:t\in\R\}$. A supplied proof $R=I$ handles the identity case.
\item \textbf{Approximate Euler for arbitrary real input.} Given $n$, one can
construct an exactly unit axis $u^{[n]}$ and a circle pair $(c^{[n]},s^{[n]})$
whose Rodrigues matrix $P_n$ satisfies
$\norm{R-P_n}_F\le2^{-n}$ and
$\norm{Ru^{[n]}-u^{[n]}}\le2^{-n}$. The axes need not form a Cauchy sequence.
\item \textbf{An exact all-input representation.} Every represented real
$R$ satisfying \eqref{eq:so} has primitive-recursively constructed unit
quaternion data $(a,v)$ satisfying
\begin{equation}\label{eq:quaternion}
 a^2+v\cdot v=1,\qquad
 R=(a^2-v\cdot v)I+2vv^T+2a\crossmat v.
\end{equation}
The vector part $v$ is allowed to vanish.
\end{enumerate}
For represented inputs, each assertion is a schema for each fixed
primitive-recursive input approximation program, uniformly in its numerical
parameters. The rational-input assertion is a single ordinary arithmetic
construction on finite codes.
\end{theorem}

\begin{theorem}[Boundary of the requested exact-axis theorem]\label{thm:boundary}
There is a fixed primitive-recursive family of rational approximation
programs for matrices $R_e\in\SO(3)$ for which no total computable procedure
can select an exact unit fixed axis uniformly in $e$. More generally,
exact unit-axis existence for all represented real rotations implies LLPO.
The resulting uniform first-order arithmetic disjunction is not provable
in $\ipra$.
\end{theorem}

Section~\ref{sec:obstruction} proves the boundary theorem in detail. It does
not assert that a particular computable rotation lacks every computable
axis. Classically each individual nonidentity matrix has a finite
nonidentity witness that can be built into an axis program. It is the
\emph{uniform constructive theorem} that fails.

\section{Arithmetic meaning of the real notation}\label{sec:arithmetic}

\subsection{The base theory and the schema convention}
We use first-order intuitionistic logic, equality, a function symbol for each
primitive-recursive definition, its defining equations, successor separation
and injectivity, and induction for quantifier-free formulas. This is the
presentation implemented as \texttt{PRAAx} in the inspected
\texttt{ConciseIPRA/PR.lean} file \cite{proveit}. Bounded quantifiers over
primitive-recursive predicates can be replaced by Boolean folds, so a
bounded finite search or a finite conjunction does not require a stronger
induction schema.

Integers are encoded by signed natural numbers; a rational is an integer
numerator and a positive denominator. Finite tuples and lists are encoded
by standard primitive-recursive pairing. Equality and order on rationals
are decided by cross multiplication. Addition, multiplication, powers,
factorials, maximum, and bounded search are primitive recursive. All
definitions can be totalized by returning a fixed default on malformed
codes or failed denominator guards. Their correctness theorems have the
appropriate well-formedness hypotheses.

Fix a primitive-recursive rational program $A(\mathbf p,k)$, where
$\mathbf p$ is a tuple of numerical parameters. A convenient regular-name
condition is
\begin{equation}\label{eq:regular}
 \abs{A(\mathbf p,k)-A(\mathbf p,l)}\le2^{-k}+2^{-l}
 \quad\text{for all }k,l.
\end{equation}
In the usual interpretation, it names a real $a$ with
$\abs{A(\mathbf p,k)-a}\le2^{-k}$. Internally, no limit object is needed:
equality, order, and algebraic equations mean their rational approximation
predicates at every precision. For example equality of regular names $A,B$
can be written $\forall k\,\abs{A(k)-B(k)}\le2^{1-k}$; nonnegative order
is expressed with the corresponding tolerances. Equivalent shifts of
precision give the same represented relation.

Every input family is fixed before forming its arithmetic formula. The
output approximation programs are explicitly composed from that fixed
program. We do \emph{not} postulate a total primitive-recursive evaluator
for arbitrary primitive-recursive program codes: diagonalization rules out
such an evaluator. A formal Lean development may manipulate PR syntax in
its metatheory, but its semantic evaluator is not thereby an object-language
PR function. This same schema distinction is made in the inspected FTA and
triangle-formula drafts \cite{drafts}. We reuse this representation
discipline, not any unverified geometric theorem from those drafts.

\subsection{Quantitative arithmetic on names}
Here are elementary constructions sufficient for all the algebraic parts.
They also make the resource bounds explicit.

\begin{lemma}[Name operations]\label{lem:operations}
The following operations have primitive-recursive rational approximations
with a prescribed error $2^{-k}$: sums and products of bounded names;
reciprocal and division given a positive rational lower bound for the
absolute denominator; and square root of a nonnegative bounded name.
\end{lemma}
\begin{proof}
For sums, request one additional bit per binary addition. If $|x|,|y|\le B$
and rational samples $a,b$ have errors at most $\eta\le1$, then
\[
 |ab-xy|\le(2B+1)\eta.
\]
An explicit precision shift makes this at most the requested error.
For a quotient, assume $|y|\ge L>0$, $|x|\le B$, and samples satisfy
$|a-x|,|b-y|\le\eta\le\min(1,L/2)$. Then $|b|\ge L/2$ and
\begin{equation}\label{eq:division-error}
 \left|\frac ab-\frac xy\right|
 \le\frac\eta L+\frac{2(B+1)\eta}{L^2}.
\end{equation}
All shifts can be found by a bounded scan: for any positive rational
$r$, an integer explicitly larger than its numerator and denominator is
enough to bound the number of doublings needed to make $2^{-m}\le r$.
There is no unbounded search for a nonzero denominator.

For square root at precision $k$, take a rational sample $b$ of $x\ge0$
with error at most $2^{-2k-4}$, replace $b$ by $\max(b,0)$, and set
$N=2^{k+2}$. Find the greatest integer $j$ in a fixed bounded range with
$j^2\le N^2b$. If $0\le x\le B$, the range
$0\le j\le (\lceil B\rceil+2)N$ suffices. The result $j/N$ differs from
$\sqrt b$ by less than $1/N$. The elementary inequality
\[
 |\sqrt x-\sqrt b|\le\sqrt{|x-b|}
\]
follows by squaring after ordering the two nonnegative rational
approximants. Hence the total error is at most $2^{-k-1}$.
This describes the square-root name directly. Alternatively the defining
property is that its nonnegative approximants have squares converging to
$x$ at the indicated rate. Uniqueness follows from
$|r-t|^2\le|r^2-t^2|$ for $r,t\ge0$.

The integer scan has a quantifier-free invariant: the current answer lies
within the processed range, its square is admissible, and every processed
larger index is inadmissible. The last finite conjunction is a Boolean
fold. Its proof uses only quantifier-free induction.
\end{proof}

The usual sum-of-squares inference is constructive: from
$x_1^2+\cdots+x_m^2=0$ one obtains each $x_i=0$ by applying the finite
inequality $x_i^2\le\sum_j x_j^2$ at an error smaller than the square of
the target precision. No trichotomy of real equality is involved.

A statement such as $x\ne0$ will not be used as a division certificate.
We use \emph{apartness}: a supplied positive rational $L$ with $|x|\ge L$.
In particular, $R\ne I$ and an explicit nonidentity bound are different
hypotheses in intuitionistic arithmetic.

\section{The finite vector and matrix identities}\label{sec:finite}

Coordinates are numbered $1,2,3$, and vectors are columns. Set
\[
 u\cdot x=u_1x_1+u_2x_2+u_3x_3,\qquad
 u\times x=(u_2x_3-u_3x_2,\ u_3x_1-u_1x_3,\ u_1x_2-u_2x_1)^T.
\]
Thus
\begin{equation}\label{eq:crossmatrix}
 \crossmat u=\begin{pmatrix}0&-u_3&u_2\\u_3&0&-u_1\\-u_2&u_1&0\end{pmatrix},
 \qquad \crossmat u x=u\times x.
\end{equation}
Expanding the coordinates gives
\begin{align}
 \crossmat u^T&=-\crossmat u,&
 \crossmat u u&=0,\label{eq:k-basic}\\
 u\times(v\times x)&=v(u\cdot x)-x(u\cdot v),\label{eq:triple}\\
 \crossmat u^2&=uu^T-(u\cdot u)I,\label{eq:k-square}\\
 \norm{u\times x}^2&=(u\cdot u)(x\cdot x)-(u\cdot x)^2.\label{eq:lagrange}
\end{align}
For example, the first coordinate of \eqref{eq:triple} expands to
$u_2(v_1x_2-v_2x_1)-u_3(v_3x_1-v_1x_3)$, which is
$v_1(u\cdot x)-x_1(u\cdot v)$ after cancelling $u_1v_1x_1$.
The other coordinates are cyclic permutations. Equation~\eqref{eq:lagrange}
is the sum of the three squared $2\times2$ minors. These are polynomial
identities over the integers, not uses of a spectral theorem.

For a $3\times3$ matrix $B$, its determinant is the signed six-term sum and
its adjugate is the transpose of the matrix of signed $2\times2$ minors.
Expanding those definitions proves
\begin{align}
 B\adj B=\adj B B&=(\det B)I,\label{eq:adjproduct}\\
 \adj(I-B)&=(1-\tr B)I+B+\adj B,\label{eq:adjshift}\\
 \det(I-B)&=1-\tr B+\tr(\adj B)-\det B,\label{eq:detshift}\\
 \adj(\adj B)&=(\det B)B.\label{eq:adjadj}
\end{align}
One way to check \eqref{eq:adjshift} is to examine a diagonal cofactor,
$(1-b_{jj})(1-b_{kk})-b_{jk}b_{kj}$: it is
$1-\tr B+b_{ii}+(\adj B)_{ii}$. An off-diagonal cofactor gives
$b_{ij}+(\adj B)_{ij}$. The determinant identity follows by collecting
the degrees $0,1,2,3$ in the six-term expansion. The two product identities
follow by the same expansions, with repeated-column terms cancelling.
All four identities also appear in the exact symbolic supplement.

\section{What orthogonality already proves without choosing an axis}
\label{sec:universal}
Suppose \eqref{eq:so}. Define
\begin{equation}\label{eq:definitions}
\begin{aligned}
 t&=\tr R,&c&=\frac{t-1}{2},&d&=1-c=\frac{3-t}{2},\\
 S&=\frac{R+R^T}{2}-cI,&T&=\frac{R-R^T}{2}=\crossmat w,\\
 w&=\frac12(R_{32}-R_{23},\ R_{13}-R_{31},\ R_{21}-R_{12})^T.
\end{aligned}
\end{equation}
In particular $R=cI+S+T$, $S$ is symmetric, and $T$ is skew-symmetric.

\begin{lemma}[Universal axis identities]\label{lem:universal}
Without any nonidentity test,
\begin{gather}
 \det(I-R)=0,\qquad \adj(I-R)=2S,\label{eq:axis-adj}\\
 RS=R^TS=S,\quad S^2=dS,\quad\tr S=d,\quad TS=ST=0,\quad\adj S=0,
 \label{eq:s-identities}\\
 \norm{R-I}_F^2=4d,\qquad c^2+w\cdot w=1,\qquad ww^T=(1+c)S.
 \label{eq:scalar-identities}
\end{gather}
Consequently $-1\le c\le1$ and $0\le d\le2$.
\end{lemma}
\begin{proof}
By \eqref{eq:adjproduct}, $R\adj R=I$. Multiplying by $R^T$ on the left
and using $R^TR=I$ gives $\adj R=R^T$, and hence also $RR^T=I$.
Equations~\eqref{eq:adjshift} and \eqref{eq:detshift} give
$\adj(I-R)=(1-t)I+R+R^T=2S$ and
$\det(I-R)=1-t+t-1=0$.
The adjugate product now gives $RS=SR=S$; multiplication by $R^T$
gives $R^TS=S$. Therefore
\[
 S^2=\left(\frac{R+R^T}{2}-cI\right)S=(1-c)S=dS.
\]
The trace identity is $t-3c=d$. The two products with $T$ vanish.
Finally \eqref{eq:adjadj} gives $\adj(2S)=0$, and since each cofactor
is homogeneous of degree two, $4\adj S=0$.

Expanding the nine squares yields
$\norm{R-I}_F^2=\tr(R^TR)-2\tr R+3=6-2t=4d$.
The entrywise inner product of a symmetric matrix with a skew-symmetric
matrix is zero, by pairing entries $(i,j)$ and $(j,i)$. Hence
\begin{align*}
 3=\norm R_F^2
 &=3c^2+2c\tr S+\tr(S^2)+2w\cdot w\\
 &=3c^2+2cd+d^2+2w\cdot w
 =1+2c^2+2w\cdot w.
\end{align*}
This proves the scalar circle identity and the stated bounds.
Using $ST=TS=0$ and \eqref{eq:k-square} in $RR^T=I$ gives
\[
 I=(c^2+w\cdot w)I+(2c+d)S-ww^T
   =I+(1+c)S-ww^T,
\]
which proves the final identity.
\end{proof}

Every column of $S$ is already an exact fixed vector. The issue at the
identity is that all three columns then vanish. Determinant zero by itself
does not construct a nonzero vector in a kernel over represented reals.
The quantitative construction below supplies precisely the missing step.

\section{Exact Euler with a supplied gap, including half-turns}
\label{sec:gap}

\begin{theorem}[Explicit unit-axis extraction]\label{thm:gap}
Suppose $d\ge\delta>0$, where $\delta$ is a supplied rational.
Approximate $S_{11},S_{22},S_{33}$ by rationals $a_1,a_2,a_3$ with error
at most $\delta/12$. Choose the least index $j$ maximizing these three
rationals, and put
\begin{equation}\label{eq:axis-construction}
 p=S_{jj},\qquad v=Se_j,\qquad q=d p,\qquad
 u=\frac{v}{\sqrt q},\qquad s=w\cdot u.
\end{equation}
Then $p\ge\delta/6$, $q\ge\delta^2/6$, and
\begin{equation}\label{eq:exactnormal}
 u\cdot u=1,\qquad Ru=u,\qquad
 c^2+s^2=1,\qquad R=cI+(1-c)uu^T+s\crossmat u.
\end{equation}
\end{theorem}
\begin{proof}
For every $i$, maximality of the rational sample gives
$S_{ii}\le S_{jj}+\delta/6$. Summing and using $\tr S=d\ge\delta$
gives $\delta\le3p+\delta/2$, so $p\ge\delta/6$.
Symmetry and $S^2=dS$ give
\[
 v\cdot v=(S^2)_{jj}=dp=q\ge\delta^2/6.
\]
Thus $\sqrt q$ has a positive numerical lower bound (for instance
$\delta/3$), and Lemma~\ref{lem:operations} applies. It follows that
$u\cdot u=1$, and $Ru=u$ follows from $RS=S$.

Since every $2\times2$ minor of $S$ vanishes, for all $i,k$,
\[
 S_{ik}S_{jj}=S_{ij}S_{jk}=v_i v_k.
\]
The cases of repeated row or column are the same equality trivially;
the other cases are cofactors. Division by $p>0$ now proves
\begin{equation}\label{eq:projector}
 S=\frac{vv^T}{p}=d uu^T.
\end{equation}
Since $TS=0$, $Tu=0$, or $w\times u=0$. The vector triple-product identity
therefore gives
\[
 0=u\times(w\times u)=w-(u\cdot w)u,
\]
so $w=su$. Substituting in $R=cI+S+\crossmat w$ gives its normal form.
Finally $c^2+w\cdot w=1$ and $u\cdot u=1$ imply $c^2+s^2=1$.
\end{proof}

\begin{corollary}[The axis is the whole fixed line]\label{cor:fixedline}
Under the gap hypothesis, $Rx=x$ if and only if $x=(u\cdot x)u$.
\end{corollary}
\begin{proof}
If $Rx=x$, multiplication by $R^T$ gives $R^Tx=x$ and hence $Sx=dx$.
Equation~\eqref{eq:projector} and $d>0$ give the conclusion. The converse
follows from $Ru=u$. Equivalently,
$(R-I)^T(R-I)=2d(I-uu^T)$.
\end{proof}

For a half-turn, $c=-1$, $d=2$, and $w=s=0$; the construction uses
$S=(R+I)$ and remains valid. A method that normalizes $w$ alone would fail
here. For the identity, $c=1$, $d=0$, and $S=w=0$; use $u=e_1$, $s=0$
when a proof of identity is supplied. There is no implicit case split
between these possibilities for arbitrary real inputs.

A particularly concrete gap certificate is
$|R_{ab}-I_{ab}|\ge\eta>0$. Equation~\eqref{eq:scalar-identities} gives
$d\ge\eta^2/4$. If $\eta$ is dyadic, so is this displayed lower bound.
The selected index $j$ is computed once from the coarse samples and then
held fixed at all output precisions. Reselecting it at each precision
would not by itself give a coherent name for one oriented axis.

\section{The unconditional exact theorem for rational matrices}
\label{sec:rational}

For rational $R$, all nine identity comparisons are decidable. If $R=I$,
take $v=e_1$, $q=1$, $c=1$, $d=b=0$. Otherwise the rational identity
$4d=\sum_{i,j}(R_{ij}-I_{ij})^2$ proves $d>0$. Choose $j$ maximizing
the \emph{exact} rational diagonal entries of $S$, so $p=S_{jj}\ge d/3$,
and set
\begin{equation}\label{eq:rationalcertificate}
 v=Se_j,\quad q=v\cdot v=dp>0,\quad b=w\cdot v.
\end{equation}
Theorem~\ref{thm:gap} now gives the entirely rational certificate
\begin{equation}\label{eq:ratcert}
 Rv=v,\qquad c^2+\frac{b^2}{q}=1,\qquad
 R=cI+\frac{d}{q}vv^T+\frac{b}{q}\crossmat v.
\end{equation}
With $u=v/\sqrt q$ and $s=b/\sqrt q$, this is exactly Rodrigues' formula.
The axis and sine coordinates lie in one quadratic extension of $\Q$;
$c$ is rational. An explicit name for the positive square root avoids
introducing an abstract algebraic closure.

All loops in the rational extraction have fixed length three or nine.
For the normalization, write $q=A/B>0$ with positive integers $A,B$.
Then $\sqrt q\ge1/B$. At requested precision $k$, set
\[
 m=k+\operatorname{bitlength}(2B),\qquad
 j=\left\lfloor\sqrt{\left\lfloor A2^{2m}/B\right\rfloor}\right\rfloor,
 \qquad r=j2^{-m}.
\]
The integer square root is obtained by a bounded scan or binary search.
It satisfies $r^2\le q<(r+2^{-m})^2$ and $r\ge1/(2B)$.
Since $|v_i|,|b|\le\sqrt q$, replacing $\sqrt q$ by $r$ in the
normalized coordinates makes an error at most
$2B2^{-m}<2^{-k}$. The exact implementation checks this interval using
rational squares. The rational approximants themselves are not claimed
to have exactly unit norm.

\paragraph{Example: a quarter-turn.}
For
\[
 R=\begin{pmatrix}0&-1&0\\1&0&0\\0&0&1\end{pmatrix}
\]
we get $c=0$, $d=1$, $S=\operatorname{diag}(0,0,1)$, $w=e_3$.
Thus $u=e_3$, $s=1$, and $Rx=u\times x+(u\cdot x)u$.

\paragraph{Example: a non-coordinate half-turn.}
Let $u=(1,2,2)^T/3$. Then
\[
 R=2uu^T-I=\frac19
 \begin{pmatrix}-7&4&4\\4&-1&8\\4&8&-1\end{pmatrix}.
\]
Here $S=2uu^T$, $c=-1$, and $w=0$. The second column is selected
by the least-index tie rule; normalizing it returns $u$. The Rodrigues
formula is $Rx=-x+2(u\cdot x)u$.

\section{Rodrigues' formula: full algebraic and geometric proof}
\label{sec:rodrigues}
Let $u\cdot u=1$ and $c^2+s^2=1$. Put $P=uu^T$, $Q=I-P$,
$K=\crossmat u$. Equations~\eqref{eq:k-basic}--\eqref{eq:k-square} give
\[
 P^2=P,\ Q^2=Q,\ PQ=QP=0,\ KP=PK=0,\ KQ=QK=K,\ K^2=-Q.
\]
The Rodrigues matrix has the useful equivalent forms
\begin{equation}\label{eq:rodrigues-forms}
 \mathcal R=P+cQ+sK
 =cI+(1-c)uu^T+sK
 =I+sK+(1-c)K^2.
\end{equation}
It fixes $u$. Since $K^T=-K$,
\[
 \mathcal R^T\mathcal R
 =(P+cQ-sK)(P+cQ+sK)
 =P+(c^2+s^2)Q=I.
\]
To check the determinant without a sign choice, direct expansion gives,
for arbitrary scalars $c,h,s$ and $n=u\cdot u$,
\begin{equation}\label{eq:det-rodrigues}
 \det(cI+h uu^T+sK)=(c+hn)(c^2+s^2n).
\end{equation}
For completeness, expand the determinant of $A+h uu^T$ by columns.
Terms using two columns from $h uu^T$ vanish because those columns are
proportional. The result is $\det A+h u^T\adj(A)u$. For $A=cI+sK$,
the six determinant terms give $\det A=c(c^2+s^2n)$, while its cofactors
give
$\adj A=c^2I-csK+s^2uu^T$. Thus
$u^T\adj A u=n(c^2+s^2n)$, proving \eqref{eq:det-rodrigues}.
Set $n=1$ and $h=1-c$ to obtain determinant one.

Every vector decomposes, without choosing a basis, as
$x=(u\cdot x)u+x_\perp$, where $u\cdot x_\perp=0$.
On this perpendicular plane, $K^2=-I$, $K$ preserves length, and
$x_\perp\cdot Kx_\perp=0$. Therefore
\[
 \mathcal R x=(u\cdot x)u+cx_\perp+s(u\times x_\perp),
\]
which is the planar rotation with oriented cosine $c$ and sine $s$.
Substituting $x_\perp=x-(u\cdot x)u$ is precisely
\eqref{eq:rodmain}. This proves the geometric meaning as well as the
matrix identity. It remains valid for $x_\perp=0$, without a division
by its length.

The same multiplication table shows
\[
 \mathcal R(u,c,s)^{-1}=\mathcal R(u,c,-s),\qquad
 \mathcal R(u,c,s)\mathcal R(u,c',s')
 =\mathcal R(u,cc'-ss',sc'+cs').
\]
These identities explain the angle addition law. The next section
connects the circle-pair formulation with ordinary radian angles by
explicit rational estimates.

% The trigonometric bridge is inserted here during assembly.
\section{An explicit radian-angle bridge using finite sums}
\label{sec:finite-trig}

The circle parameters $c,s$ in the algebraic theorem can be replaced by a
real angle $\theta$ with $\cos\theta=c$ and $\sin\theta=s$.  Here is a
construction with primitive recursive precision bounds.  It does not invoke
the fundamental theorem of calculus.  Absolute values of complex numbers
are shorthand for the elementary Euclidean norm; all the estimates below
can equivalently be checked after squaring and clearing positive rational
denominators.

\subsection{Sine and cosine from finite Taylor polynomials}

Put
\[
 T_m(t)=\sum_{k=0}^{m}\frac{(\mathrm i t)^k}{k!}.
\]
For an integer $B\geq1$, $|t|\leq B$, and $m\geq2B$, the tail bound is
\[
 \sum_{k=m+1}^{\infty}\frac{B^k}{k!}
 \leq 2\frac{B^{m+1}}{(m+1)!}
 \leq 2^{\,2B^2+2B-m}.                                      \tag{T1}
\]
Indeed, after index $2B$ successive terms have ratio at most $1/2$,
$B^{2B}\leq2^{2B^2}$, and a finite geometric-sum bound gives the displayed
inequality before taking any limit.  Thus
$m=2B^2+2B+n+4$ already makes the rational tail at most $2^{-n-4}$.
The resulting rapid Cauchy name defines
\[
 E(t)=\cos t+\mathrm i\sin t.
\]

For completeness, this definition also gives a uniform algorithm when
$t$ is supplied by a real name, not a rational number.  Suppose
$|t|\leq M$, where $M\geq1$ is an integer, and set $B=M+1$.
Evaluate the input name at precision
$n+2B^2+2B+8$, obtaining a rational $q$; evaluate $T_m(q)$ with
$m=2B^2+2B+n+6$.  The error in a polynomial value caused by changing
$t$ to $q$ is at most
\[
 |t-q|\sum_{j\geq0}\frac{B^j}{j!}
 \leq |t-q|\,2^{2B^2+2B+2}.
\]
This and (T1) give an error smaller than $2^{-n}$.  Factorials, finite
sums, and the displayed precision transformers are primitive recursive.

We record the elementary identities needed below.  The finite Cauchy
product and the binomial theorem give
\[
 E(t+u)=E(t)E(u),\qquad E(-t)=\overline{E(t)},\qquad E(0)=1.   \tag{T2}
\]
There is no appeal to a theorem about infinite series hidden here: if
$|t|,|u|\leq B$, all terms of total degree at most $m$ agree exactly in
$T_m(t)T_m(u)$ and $T_m(t+u)$.  The absolute sum of the remaining terms is
at most the tail in (T1) with $2B$ in place of $B$.  For example,
$m=8B^2+4B+n+4$ suffices: using the preceding bound on the exponential
sum, the errors from replacing each factor by its Taylor polynomial,
the finite Cauchy-product remainder, and the tail for $t+u$ together
are at most $3\cdot2^{-n-4}$.  This proves (T2) at any requested
precision.  Consequently
$E(t)\overline{E(t)}=1$, so $|E(t)|=1$.

For $|a|\leq1$, the same finite-series estimates give
\[
 |E(a)-1|\leq 2|a|,\qquad
 |E(a)-1-\mathrm ia|\leq |a|^2,                             \tag{T3}
\]
because $\sum_{k\geq1}1/k!\leq2$ and
$\sum_{k\geq2}1/k!\leq1$.  Factoring $E(h)$ as a product of $K$ copies
of $E(h/K)$, with a supplied integer $K\geq\max(1,|h|)$, and telescoping
the product yields the convenient uniform bound
\[
 |E(a)-E(b)|\leq 3|a-b|.                                   \tag{T4}
\]
The harmless constant $3$ leaves slack for rational approximations.

\subsection{Arctangent from dyadic rational sums}

For a positive integer $N$, define
\[
 A_N(t)=\frac{t}{N}\sum_{j=0}^{N-1}
                  \frac{1}{1+(tj/N)^2}.                    \tag{T5}
\]
These are finite rational expressions for rational $t$; all denominators
are at least one.  Suppose $|t|\leq M$, $M\geq1$.  The exact rational
identity for $f(x)=1/(1+x^2)$ gives, for $|x|,|y|\leq M$,
\[
 |f(x)-f(y)|
 =\frac{|x-y|\,|x+y|}{(1+x^2)(1+y^2)}
 \leq2M|x-y|.
\]
Write $h=t/N$.  Their difference is exactly
\[
 A_N(t)-A_{2N}(t)
 =\frac h2\sum_{j<N}\bigl(f(jh)-f(jh+h/2)\bigr).
\]
The preceding bound therefore gives
\[
 |A_N(t)-A_{2N}(t)|\leq\frac{M^3}{2N}.
\]
Summing the finite geometric progression proves that the dyadic sums
define a rapid Cauchy real $A(t)$, with
\[
 |A(t)-A_N(t)|\leq\frac{M^3}{N}
 \quad\hbox{when $N$ is a power of two}.                    \tag{T6}
\]
The notation $A$ is the usual arctangent: (T5) is also a direct
constructive definition by rational integration.  No integration theorem
is required for the identities used here.

Input error is especially easy to control.  For $0\leq r\leq1$, put
$g(t,r)=t/(1+t^2r^2)$.  Direct subtraction gives
\[
 g(t,r)-g(u,r)
 =\frac{(t-u)(1-tu r^2)}{(1+t^2r^2)(1+u^2r^2)},
\]
and $(t^2+u^2)r^2\geq2|tu|r^2$ shows that its absolute value is at most
$|t-u|$.  Hence every $A_N$, and then $A$, is $1$-Lipschitz.
For a real input satisfying $|t|\leq M$, obtain a rational $q$ at
precision $n+3$, set $B=M+1$, and use the dyadic mesh
\[
 N=2^{n+3B+4}.
\]
Since $B^3\leq2^{3B}$, (T6) and the input-error estimate give an
approximation to $A(t)$ with error less than $2^{-n}$.  This is an
explicit primitive recursive name transformer.

\subsection{The Cayley identity without differentiation}

Define the unit complex rational function
\[
 w(x)=\frac{1+\mathrm ix}{1-\mathrm ix}
     =\frac{1-x^2+2\mathrm ix}{1+x^2}.
\]
We prove
\[
                     E(2A(t))=w(t).                        \tag{T7}
\]
Fix a dyadic $N\geq2M$, put $h=t/N$, and $x_j=jh$.
Because $x_jh\geq0$, the number
$D_j=1+x_j^2+x_jh$ is at least one.  Exact multiplication gives
\[
 \frac{w(x_j+h)}{w(x_j)}
       =\frac{D_j+\mathrm ih}{D_j-\mathrm ih}.
\]
Writing $r_j=h/D_j$, the difference of the right side from
$1+2\mathrm ir_j$ has norm at most $2h^2$.
Also
\[
 \left|r_j-\frac{h}{1+x_j^2}\right|
 =\frac{|x_j|h^2}{D_j(1+x_j^2)}\leq Mh^2.
\]
Thus its difference from $1+2\mathrm i h/(1+x_j^2)$ is at most
$2(M+1)h^2$.  Since $|2h/(1+x_j^2)|\leq1$, (T3) shows that
\[
 \left|\frac{w(x_j+h)}{w(x_j)}
       -E\!\left(\frac{2h}{1+x_j^2}\right)\right|
 \leq(2M+6)h^2.                                            \tag{T8}
\]

All factors in the two products have norm one.  Telescoping the products
in (T8), using $w(0)=1$ and (T2), gives
\[
 |w(t)-E(2A_N(t))|
 \leq \frac{(2M+6)M^2}{N}.
\]
By (T4) and (T6),
\[
 |w(t)-E(2A(t))|
 \leq\frac{8M^3+6M^2}{N}
 \leq\frac{16(M+1)^3}{N}.                                  \tag{T9}
\]
For example, $N=2^{n+3(M+1)+5}$ makes this error at most $2^{-n-1}$
and ensures $N\geq2M$.  In the later uniform case $M=2$, the smaller
choice $N=2^{n+9}$ already makes it less than $2^{-n}$.
Thus an explicit finite estimate at any requested precision proves (T7).
This is a finite error argument throughout, not an application of the
fundamental theorem of calculus or an existence theorem for differential
equations.

Define $\pi=4A(1)$, the standard arctangent normalization of the radian
constant.  Equation (T7) gives $E(2A(1))=\mathrm i$, so (T2) gives
\[
                           E(\pi)=-1.                       \tag{T10}
\]

\subsection{A finite chart choice produces an angle}

Let $c^2+s^2=1$.  Compute once a rational $q$ with
$|q-c|\leq1/4$, and make the decidable rational comparison $q\geq0$.
If this holds, $1+c\geq3/4$; set
\[
             t=\frac{s}{1+c},\qquad \theta=2A(t).
\]
If $q<0$, then $1-c>3/4$; set
\[
             t=\frac{s}{1-c},\qquad \theta=\pi-2A(t).
\]
In either case $|t|\leq4/3$, so one may use the uniform bound $M=2$
in all computations.  In the first chart, $w(t)=c+\mathrm is$; in the
second, $-\overline{w(t)}=c+\mathrm is$.  These follow by multiplying
through by $(1+c)^2+s^2=2(1+c)$ or
$(1-c)^2+s^2=2(1-c)$, respectively.  Equations (T2), (T7), and (T10)
therefore give $E(\theta)=c+\mathrm is$ in both charts.

This is a valid existential angle constructor.  Its selected angle may
depend on the supplied name and need not be a globally continuous
single-valued principal argument on the circle.  Both chart formulas
produce a correct angle wherever that chart is selected.  The chart is
chosen once at a fixed input precision and retained at every later
output precision.

\paragraph{Arithmetic audit.}
Every sum, product, dyadic precision, and number of iterations appearing
above has a displayed primitive recursive bound.  The finite binomial,
geometric-sum, and telescoping identities use quantifier-free induction
after rational denominators have been cleared.  The output precision is
a free natural parameter in those induction arguments.  Universal
quantification over precision is introduced after each finite estimate
has been proved; induction on a Cauchy statement containing such a
quantifier is not used.


\section{An unconditional exact quaternion representation}
\label{sec:quaternion}

This section gives a stronger all-input exact conclusion than determinant
zero, while respecting the axis obstruction. Approximate $c$ by a rational
$z$ with $|z-c|\le1/4$, and perform the decidable rational test $z\ge0$.

If $z\ge0$, then $1+c\ge3/4$. Set
\[
 a=\sqrt{(1+c)/2},\qquad v=w/(2a).
\]
Every denominator is bounded away from zero. From
$w\cdot w=1-c^2$ and $ww^T=(1+c)S$,
\[
 v\cdot v=(1-c)/2,\qquad 2vv^T=S,\qquad 2av=w.
\]
Hence $a^2+v\cdot v=1$, and the right side of
\eqref{eq:quaternion} is $cI+S+\crossmat w=R$.

If $z<0$, then $c<1/4$, so $d>3/4$. Apply the exact-axis construction
with, for example, $\delta=1/2$, obtaining $u,s$. Set
\[
 v=\sqrt{d/2}\,u,\qquad a=s/\sqrt{2d}.
\]
The identity $s^2=1-c^2=d(1+c)$ gives $a^2=(1+c)/2$.
Again $v\cdot v=d/2$, $2vv^T=S$, and $2av=su=w$, proving
\eqref{eq:quaternion}.

Both branches are primitive recursive on names and use fixed positive
denominator bounds. This is a construction of a representative quaternion;
it is not a claim of a globally continuous, extensional sign choice on
$\SO(3)$. Different sufficiently close input names may select different
valid signs. At the identity the first branch gives $(1,0)$.
Normalizing the vector part there would reintroduce the forbidden exact
axis selection. This is why the quaternion theorem and
Theorem~\ref{thm:boundary} are compatible.

\section{Uniform finite-precision Euler theorem}
\label{sec:approximate}

\begin{theorem}[A quantitative approximation without a gap promise]
For every integer $n\ge0$, let $\eps=2^{-n}$. Obtain a rational $a$ with
$|a-d|\le\eps^2/32$ and decide whether $a\le\eps^2/8$.
\begin{enumerate}[label=(\alph*)]
\item If yes, return $u^{[n]}=e_1$, $c^{[n]}=1$, $s^{[n]}=0$ and $P_n=I$.
\item If no, apply Theorem~\ref{thm:gap} with
$\delta=\eps^2/16=2^{-(2n+4)}$, and return its exact data and $P_n=R$.
\end{enumerate}
Then $u^{[n]}$ is exactly unit, $(c^{[n]})^2+(s^{[n]})^2=1$, and
\[
 \norm{R-P_n}_F\le\eps,\qquad
 \norm{Ru^{[n]}-u^{[n]}}\le\eps.
\]
\end{theorem}
\begin{proof}
In the first branch $d\le5\eps^2/32$, and therefore
\[
 \norm{R-I}_F^2=4d\le5\eps^2/8<\eps^2.
\]
The residual bound follows either from Cauchy--Schwarz applied row by row
or by observing that $(R-I)e_1$ is one column of $R-I$.
In the second branch $d>3\eps^2/32\ge\eps^2/16$, so the exact theorem
applies. Its residual and reconstruction errors are zero.
\end{proof}

The branch tolerance is quadratic in the requested matrix error. The
two branches overlap in a safe tolerance region; no exact comparison of
a real with zero occurs. For every fixed $n$, the resulting axis is
represented by one coherent PR name. Across different $n$, neither the
chosen direction nor its orientation is required to converge. Such a
requirement would be a different theorem.

\subsection{A completely finite rational output version}
In the large branch, approximate the exact $u,c,s$ coordinatewise with
error $\eta=2^{-n-5}$, obtaining rational $U,C,V$. Since
$|u_i|,|c|,|s|\le1$ and $\eta\le1$, elementary expansion gives
\begin{align}
 |U\cdot U-1|&\le9\eta,&|C^2+V^2-1|&\le6\eta,\label{eq:finite-unit}\\
 \max_i|(RU-U)_i|&\le4\eta,\label{eq:finite-resid}\\
 \max_{i,j}|\mathcal R(U,C,V)_{ij}-R_{ij}|&\le14\eta.\label{eq:finite-rod}
\end{align}
Here $\mathcal R(U,C,V)$ means the polynomial expression, even if the
rational samples are not exactly normalized. To verify the constant
$14$, the $C I$ term costs at most $\eta$; the term
$(1-C)U_iU_j$ costs at most $3(3\eta)+\eta=10\eta$; and the
skew term costs at most $3\eta$. The diagonal cross term is zero,
so the same bound covers every entry. Norm and circle bounds use
$|a^2-b^2|\le3\eta$ when $|b|\le1$, $|a-b|\le\eta\le1$.
Every right side is at most $2^{-n}$. If the input matrix is also
replaced by rational samples of error $\eta$, the reconstruction bound
becomes $15\eta$, still at most $2^{-n}$.

In the small branch the exact rational output $(e_1,1,0)$ already has
zero normalization errors and the previously proved residual and
reconstruction bounds. Thus the approximate theorem has a purely finite
arithmetic endpoint, not only an informal claim about real limits.

\section{Why unrestricted exact axis selection is not an iPRA theorem}
\label{sec:obstruction}

\subsection{A self-contained reduction to LLPO}
Let $(\alpha_m)$ be a binary sequence containing at most one $1$.
The lesser limited principle of omniscience, LLPO, asserts the disjunction
\begin{equation}\label{eq:llpo}
 (\forall k\,\alpha_{2k}=0)\ \vee\ (\forall k\,\alpha_{2k+1}=0).
\end{equation}
This is an additional constructive principle, not a permitted real
equality case split; see \cite{bridgesrichman} for the standard terminology.
Define
\[
 x=\sum_{k=0}^\infty 2^{-k-2}\alpha_{2k},\qquad
 y=\sum_{k=0}^\infty 2^{-k-2}\alpha_{2k+1},\qquad D=1+x^2+y^2.
\]
The tails after $k<N$ are bounded by $2^{-N-1}$. If the sequence is
primitive recursive in parameters, so are these real names. In particular
$D\ge1$ gives a uniform division bound. Form
\begin{equation}\label{eq:cayley}
 R=\frac1D\begin{pmatrix}
 1+x^2-y^2&2xy&2y\\
 2xy&1-x^2+y^2&-2x\\
 -2y&2x&1-x^2-y^2
 \end{pmatrix}.
\end{equation}
For every $x,y$, this is a proper orthogonal matrix. One verification is
\eqref{eq:quaternion} with the unnormalized quaternion $(1,x,y,0)$ and
division by its squared norm $D$. Alternatively, if
$K=\crossmat{(x,y,0)^T}$, direct expansion gives
\[
 R=I+\frac{2(K+K^2)}D,\qquad (I-K)R=I+K,
\]
using $K^3=-(x^2+y^2)K$. These identities and their determinant
expansions are finite rational identities.

Suppose an exact unit fixed vector $u$ is supplied. The last identity
and $Ru=u$ imply $(I-K)u=(I+K)u$, hence $Ku=0$.
If $\alpha_{2k}=1$, the at-most-one condition gives
$x=2^{-k-2}>0$ and $y=0$, so $u_2=u_3=0$.
If $\alpha_{2k+1}=1$, it gives $y=2^{-k-2}>0$ and $x=0$, so
$u_1=u_3=0$. Each division used in these implications has the
displayed dyadic lower bound.

Take rational samples $q_i$ of the three $u_i$ with error at most $1/16$.
Some $|q_i|>1/2$. Indeed, if all three rational comparisons failed,
each $|u_i|\le9/16$, and
$u\cdot u\le243/256<1$, a contradiction. The finite comparisons
are decidable, so choose the least such $i$; then $|u_i|>7/16$.
If $i=1$, no odd-indexed $1$ is possible, so all odd terms are zero.
If $i=2$ or $3$, no even-indexed $1$ is possible, so all even terms
are zero. Decidability of each binary term turns these exclusions into
the respective universal equalities. This constructs \eqref{eq:llpo}.

This proves a lower bound on the logical strength of axis selection.
It does not claim that LLPO alone is an exact classification of every
possible arithmetization of the problem.

\subsection{A precise first-order arithmetic obstruction}
To avoid quantification over arbitrary real programs, fix an ordinary
effective enumeration $(\varphi_e)$ of partial computable unary functions.
Let $E(e,k)$ be the primitive-recursive predicate that $\varphi_e(e)$
first halts at step $k$ with output $0$, and let $O(e,k)$ be the analogous
predicate for output $1$. Only a bounded simulation is required.
Put $\alpha_{2k}=E(e,k)$ and $\alpha_{2k+1}=O(e,k)$ in the preceding
construction. At most one term is $1$, uniformly provably by the
deterministic simulation. This produces one fixed PR matrix-name family
$R_e(n)$.

The arithmetic consequence of exact axis selection for this family is
\begin{equation}\label{eq:H}
 H:\quad\forall e\bigl[(\forall k\ E(e,k)=0)\ \vee\
                          (\forall k\ O(e,k)=0)\bigr].
\end{equation}
This formula has no real variables and no universal PR evaluator.
It is not number-realizable. A realizer would compute, for each $e$, a
disjunction tag $b(e)\in\{0,1\}$ selecting a bit that
$\varphi_e(e)$ cannot output. Specifically the first disjunct forbids
output $0$, and the second forbids output $1$. The tag is a total
computable function. Let $k$ be an index for that function. Then
$\varphi_k(k)=b(k)$, contradicting the bit forbidden by its own tag.

Every theorem of intuitionistic arithmetic with the stated PR symbols
is number-realizable; in particular this holds for $\ipra$. The soundness
proof interprets a disjunction by a tag and a realizer of the selected
side, an existential by a witness and a realizer, and a universal by an
effective operation on inputs. Logical rules preserve these clauses.
The PR defining equations are true decidable equations, and induction
is realized by finite iteration. This is the usual Kleene--Nelson
realizability argument \cite{kleene,moschovakis}. Consequently $H$ is
not provable in $\ipra$.

Any exact-axis assertion in the schema interpretation of
Section~\ref{sec:arithmetic}, with uniformly constructed output names,
would prove $H$ by the finite-coordinate argument above. The same
holds for an explicit first-order coding of output names that permits
their claimed effective evaluation. Therefore such an assertion cannot
be proved. Even allowing arbitrary computable, rather than merely PR,
uniform output procedures would not help.

\subsection{Two distinctions that prevent misleading conclusions}
The obstruction is not the demand for an extensional or continuous
\emph{canonical} axis. The argument permits either orientation and any
name-dependent choice. It rules out every uniform computable choice
with the exact fixed-axis property.

Nor is it a counterexample to the classical theorem. Each individual
computable matrix has a computable axis classically: the identity has
$e_1$, and a nonidentity matrix has some finite nonidentity witness
that can be hardcoded in Theorem~\ref{thm:gap}. The unavailable operation
is finding the appropriate case or witness uniformly.

There is also a constructive double-negation statement, when axis
existence is expressed in a setting with the indicated names:
$\neg\neg\exists u\,(u\cdot u=1\wedge Ru=u)$.
If no axis existed, no entry of $R-I$ could be apart from zero, because
an apart entry yields an axis. Located rational approximation then
forces every entry to be zero. But $I$ has the axis $e_1$, a contradiction.
This supplies no exact axis program, and is not substituted for the
witness-producing theorem.

% The proof audit, Lean plan, and bibliography are inserted during assembly.
\section{Why the positive proofs stay inside iPRA}\label{sec:proof-audit}

The claim of an $\ipra$ proof is more specific than the claim that the
displayed numerical algorithms are primitive recursive. The arithmetic
must also prove their correctness with its permitted induction. We give
the reduction explicitly, so that an analytic-looking intermediate line
does not silently introduce stronger induction or a completeness axiom.

\subsection{Finite algebra certificates and approximation}
An integer polynomial identity in a fixed finite number of variables is
proved by expanding, collecting like monomials, and applying the
commutative-ring laws. These laws for encoded integers and rationals
follow by quantifier-free induction from the PR equations for addition
and multiplication. A finite list of coefficients is itself a PR object.
If one wants a uniform identity-checking procedure, polynomial addition,
multiplication and evaluation are PR finite-list operations; the proof
that evaluation respects them is a quantifier-free induction after the
finite lists are coded.

Every use of an orthogonality hypothesis in
Section~\ref{sec:universal} has an explicit polynomial residual identity.
For example, for an arbitrary matrix $A$, put
$E=A^TA-I$ and $D=\det A-1$. The exact identity
\begin{equation}\label{eq:adj-residual}
 \adj A-A^T=D A^T-E\adj A
\end{equation}
follows from $A\adj A=(\det A)I$ by multiplication with $A^T$.
Thus approximate orthogonality and determinant one imply approximate
cofactor inversion with a finite, computable error bound. The remaining
universal identities are obtained by the displayed additions and
multiplications, so the same argument applies to each one. There is no
step that converts determinant zero into a nonzero kernel vector by a
classical rank theorem.

Here is a general bound that makes this elimination of real notation
mechanical. If $P(X)=\sum_\alpha a_\alpha X^\alpha$, $B\ge1$,
$|x_i|,|y_i|\le B$, and $\max_i|x_i-y_i|\le\eta$, then
\begin{equation}\label{eq:poly-lipschitz}
 |P(x)-P(y)|\le L(P,B)\eta,\qquad
 L(P,B)=\sum_\alpha |a_\alpha|\,|\alpha|\,B^{|\alpha|-1},
\end{equation}
where constant monomials contribute zero. To prove this, write the
difference of a product of $m$ factors as a sum of $m$ products with one
difference factor. Each such term is bounded by $B^{m-1}\eta$.
Sum over the monomials. This is finite telescoping and finite summation.
Rational coefficients cause no difficulty after positive denominators
are cleared.

For products, quotients with a displayed lower bound, and square roots,
use Lemma~\ref{lem:operations}. Reading an expression tree from its root
to its inputs assigns an explicit required input precision at each
node. A product multiplies its error budget by the displayed magnitude
bound; a quotient uses \eqref{eq:division-error}; a square root requests
twice the precision plus a fixed shift. All expression trees in the
matrix proof have fixed finite size. Their resulting precision
transformers are PR in $n,h$ and the finite magnitude bounds.

For instance, after clipping the nine rotation-entry approximants to
$[-1,1]$, entry error $\eta$ gives error at most $3\eta$ in every entry
of $S$, at most $2\eta$ in $d$, and at most $\eta$ in every component
of $w$. With $\delta=2^{-h}$, sample the input at precision $h+6$.
Then $3\eta\le3\delta/64<\delta/12$, as required by the pivot proof.
Clipping is a rational min/max operation and does not increase error
when the represented entry lies in $[-1,1]$. Orthogonality proves that
entry bound by a sum of squares.

The chosen column has $q=v\cdot v=dp\ge\delta^2/6$ and
$\sqrt q\ge\delta/3$. Thus every later precision is bounded before
execution. The stored index $j$, obtained by three rational comparisons,
is independent of the requested output precision. This proves the
regularity of the output name as well as its algebraic properties.

\subsection{Where induction occurs}
The following table lists the induction tasks. ``Finite'' means a
primitive-recursive bound is part of the construction.

\begin{longtable}{@{}p{0.27\textwidth}p{0.64\textwidth}@{}}
\toprule
Task & Quantifier-free invariant or identity \\
\midrule
\endhead
Integer and rational arithmetic & Defining equations for addition,
multiplication and comparison; denominator positivity and the cleared
integer equality or inequality.\\[3pt]
Finite sums and products & Equality between the accumulator and the
finite folded expression; triangle or product error bounded by the
corresponding finite error sum.\\[3pt]
Maximum and pivot & The selected rational sample dominates every sample
already processed. There are only three axis samples.\\[3pt]
Integer square root & Current candidate is in range and has admissible
square; every larger processed candidate is inadmissible, encoded by
a bounded Boolean fold.\\[3pt]
Taylor estimates & Factorial-ratio bounds, the finite geometric-sum
bound, and the finite binomial convolution at the requested degree.\\[3pt]
Arctangent mesh comparison & The finite rational difference between
two adjacent dyadic meshes has the displayed bound.\\[3pt]
Cayley product & The difference of the two products of the first $j$
factors is bounded by the sum of their first $j$ errors.\\[3pt]
PR input simulation in the obstruction & Deterministic bounded
simulation has at most one first halting step and one output.\\
\bottomrule
\end{longtable}

When an invariant contains a bounded ``for every processed index'',
replace it with its PR Boolean fold. The induction formula is then
quantifier free in the full PR language. Its context may contain
universally quantified assumptions about the input names; the
quantifier-free restriction is on the induction formula, not on the
whole context.

For additional precision, fix a requested output precision $n$ before
performing an induction. The finitely many input samples needed by that
calculation form a coded finite table. Its regularity and polynomial
residual bounds, at the sampled indices, are decidable rational
conditions. Let $\mathsf{GoodTable}$ be their Boolean conjunction. The
actual induction can be applied to a formula of the form
\[
 \mathsf{GoodTable}=1\ \longrightarrow\
 \mathsf{AccumulatorBound}(j,n,h,\mathsf{table})=1.
\]
All quantities and loop bounds are PR. Instantiate the input's
universally quantified premises at those indices to establish
$\mathsf{GoodTable}=1$. Only after finishing the finite calculation do
we universally quantify over $n$. We never apply induction to a
statement asserting a Cauchy property at all precisions.

\subsection{The actual arithmetic conclusions}
For rational input the first endpoint is particularly simple. Encode a
rational matrix by a natural number. Let $\mathsf{RatSO3}(r)$ be the
decidable test of its nine orthogonality equations and its determinant.
Let $V(r),Q(r),C(r),D(r),B(r)$ be the totalized PR extraction functions.
The statement
\[
 \forall r\bigl(\mathsf{RatSO3}(r)\longrightarrow
   \mathsf{RatCertificate}(r,V(r),Q(r),C(r),D(r),B(r))\bigr)
\]
is an ordinary arithmetic theorem, with a quantifier-free matrix of
decoded rational equalities and inequalities. The predicate
$\mathsf{RatCertificate}$ is exactly \eqref{eq:ratcert}, together with
$Q=V\cdot V>0$ and $D=1-C$. The normalization program adds universal
finite-precision bounds; interpreting its explicit regular names gives
the exact unit-axis assertion.

For each fixed PR real-matrix family $A$, the separated result has the
form
\begin{equation}\label{eq:object-schema}
\begin{split}
 \forall\mathbf p\,\forall h\bigl(&\mathsf{Regular}(A,\mathbf p)
 \wedge\mathsf{Orthogonal}(A,\mathbf p)
 \wedge\mathsf{DetOne}(A,\mathbf p)
 \wedge\mathsf{Gap}(A,\mathbf p,h)\\
 &\longrightarrow\mathsf{UnitAxisAndRodrigues}
       (A,U_A,C_A,S_A,\mathbf p,h)\bigr).
\end{split}
\end{equation}
Each displayed predicate abbreviates its first-order rational
approximation formula; its unbounded quantifiers range over numerical
precisions. The output symbols are the concrete PR definitions already
constructed. There is no existential quantifier ranging over an
unspecified function space. The finite proof described above is a
derivation of each such schema instance. The unconditional approximation
and quaternion statements have the same form, with the gap replaced by
the robust finite branch supplying it where needed.

This explains both parts of the logical status. The positive proofs
translate into the stipulated arithmetic because every computational
loop and invariant is explicit. The unrestricted exact unit-axis claim
does not become true by calling it a schema: the fixed PR family in
Section~\ref{sec:obstruction} is already a forbidden instance.

\section{Lean formalization: existing interfaces and a completion path}
\label{sec:lean}

\subsection{Pinned repository and reusable infrastructure}
The repository was inspected at commit
\begin{center}\small\ttfamily
c1bda76f7f86d2ade505355980826f625e36fff3.
\end{center}
The root toolchain and Mathlib dependency both specify version
\texttt{v4.32.0}. The relevant source names and their roles are:
\begin{center}
\begin{tabular}{@{}p{0.40\textwidth}p{0.52\textwidth}@{}}
\toprule
Interface & Role \\
\midrule
\texttt{PR}, \texttt{PTm} & PR function symbols and first-order terms.\\
\texttt{Fm}, \texttt{Prf} & Formulas and intuitionistic derivations.\\
\texttt{PRAAx.ind}, \texttt{PRA.qfInd} & Induction with an explicit
\texttt{QF} side condition.\\
\texttt{pra\_theorem\_iff} & Transport from iPRA to the scalar theory
$T_{\rm eq}$.\\
\texttt{pra\_iff\_native} & Transport from iPRA to the native concise
calculus.\\
\bottomrule
\end{tabular}
\end{center}
These are actual inspected declarations \cite{proveit}. They transport
an existing object-language proof. They do not translate an arbitrary
theorem about Mathlib's real numbers into an iPRA derivation.

\subsection{Artifacts supplied now}
\texttt{lean/RodriguesKernel.lean} contains substantive algebraic proof
scripts over an arbitrary commutative ring. They cover dot and cross
identities, the vector triple product, Lagrange's identity, fixation of
the unit axis, preservation of dot products, and determinant one. The
scripts use explicit coordinates and \texttt{ring}. They contain no
\texttt{sorry} declaration or added axiom, but they have \emph{not} been
run through Lean in this environment.

\texttt{lean/ObjectLanguageBoundary.lean} uses the repository's actual
\texttt{PR} type to specify rational-name symbols. It also gives a small
transport lemma taking an already supplied iPRA derivation to the
native calculus. This is explicitly an interface example, not the
missing geometry derivation.

The environment had no \texttt{lean} or \texttt{lake} executable.
Accordingly no claim is made that these source files elaborate, that
their dependencies have passed a kernel audit, or that the entire
theorem is formalized. The supplied detailed plan identifies the
remaining work; this is the requested outline with a concrete path to
completion.

\subsection{A clear sequence of formalization milestones}
\begin{enumerate}
\item \textbf{Compile the finite algebra prototype.} In the pinned
ProveIt environment, compile \texttt{RodriguesKernel.lean}; resolve any
syntax or simplifier differences. Inspect \texttt{\#print axioms} for
the final declarations. This checks the algebra specification but does
not yet establish object-language derivability.
\item \textbf{Implement encoded rational arithmetic.} Construct
\texttt{PR} definitions for signed integers, positive-denominator
rationals, comparison, powers, finite folds and bounded search. Prove
their arithmetic and order properties as \texttt{Prf PRAAx} terms,
using \texttt{PRA.qfInd} for the displayed finite invariants.
\item \textbf{Turn fixed polynomial identities into derivations.}
Either write direct ring-law derivations for the finite collection here,
or implement a certificate elaborator for coefficient normalization.
Its output must be an actual object-language derivation. A Python
symbolic success message and a semantic Lean \texttt{ring} proof are
independent checks, not substitutes for that output.
\item \textbf{Finish the rational theorem first.} Its input predicate,
identity test, pivot, and certificate are all finite and rational.
Prove the totalized extractor correct, then prove the dyadic-root
normalization interval and the precision bounds of
Section~\ref{sec:rational}. This gives a compact exact initial target.
\item \textbf{Formalize regular-name arithmetic.} Encode the predicates
of Section~\ref{sec:arithmetic} and prove the explicit operation
transformers. Instantiate universal input premises at the finitely
needed precision indices. Keep target precision free throughout every
induction, and universally introduce it afterward.
\item \textbf{Formalize the gap construction.} Prove the universal
cofactor identities, the three-sample pivot bound, and the selected
column's positive norm. Compose the name operations to obtain
$U_A,C_A,S_A$, then build \eqref{eq:object-schema}. The rank-one
argument is the vanishing-minor proof given above; no rank API or
eigenvector choice is needed.
\item \textbf{Add the robust branches.} Formalize the small/large
$d$ comparison for finite-precision Euler and the two $c$ charts for
the exact quaternion theorem. Their decisions are rational and their
proofs use the displayed overlapping thresholds.
\item \textbf{Add radian angles.} Define finite Taylor polynomials and
dyadic arctangent sums as PR functions. Prove (T1)--(T9) at fixed
precision using the listed invariants; add the two circle charts. This
proves both the given-angle version of Rodrigues and reconstruction
of a radian angle from the extracted circle pair.
\item \textbf{Transport only the finished derivation.} Close the
arithmetic schema instance, construct its actual
\texttt{[] $\vdash$[PRAAx] formula} proof term, and apply
\texttt{pra\_theorem\_iff} or \texttt{pra\_iff\_native}. Audit the
induction side conditions and ensure there is no new arithmetic axiom.
\end{enumerate}

The nonprovability theorem is a separate metamathematical development:
formalize bounded universal simulation and its unique first halt,
the Cayley reduction, and the realizability soundness theorem. It is
not needed to validate any positive algorithm. Its explicit arithmetic
target is \eqref{eq:H}; the note does not pretend that the concise
repository already contains its formal nonprovability proof.

\section{Verification, scope of evidence, and reproducibility}
\label{sec:verification}

The executable certificate module uses exact rational arithmetic.
It accepts rational $R$, verifies \eqref{eq:so}, extracts
$(v,q,c,d,b)$, and checks every equality in \eqref{eq:ratcert} exactly.
The normalization routine returns rational lower and upper square-root
bounds and a rational coordinate error certificate. It never treats
a floating-point tolerance as exact equality.

The test suite includes every nonzero integer quaternion in
$\{-2,-1,0,1,2\}^4$: 624 input quaternions, including four identity
inputs and 124 half-turn inputs. There are 31 additional edge inputs,
including small-angle and near-half-turn scales through $2^{128}$.
The initial normalization suite checks 1,965 combinations of those
inputs with requested precisions $0,8,32$, and rejects three invalid
matrices. The JSON reports record the final exact counts, including
the additional finite-precision branch checks.

The symbolic supplement checks the finite adjugate, dot/cross,
Rodrigues, Cayley and quaternion numerator identities by exact polynomial
arithmetic. These checks provide useful independent evidence for the
finite calculations. They do not prove the quantified $\ipra$ theorem,
certify the analytic estimates automatically, or establish Lean
elaboration. Those limits are stated in the source files and the
formalization status file.

To reproduce the computational evidence and PDF, from the package root:
\begin{lstlisting}
python3 python/verify.py
python3 python/verify_polynomials.py
pdflatex -interaction=nonstopmode -halt-on-error euler_rodrigues_ipra.tex
pdflatex -interaction=nonstopmode -halt-on-error euler_rodrigues_ipra.tex
\end{lstlisting}
The symbolic checker needs SymPy; the rational certificate checker uses
only Python's standard library. The package includes the LaTeX sources,
compiled PDF, Lean prototype and plan, exact-arithmetic scripts,
verification reports, and a source manifest. The manifest pins the
inspected repository files and the two relevant drafts. The repository
was read, not modified.

\begin{thebibliography}{9}
\raggedright
\small
\bibitem{eulerwiki}
Wikipedia contributors, \emph{Euler's rotation theorem}, user-specified
statement, accessed November 14, 2026.
\url{https://en.wikipedia.org/wiki/Euler%27s_rotation_theorem}.

\bibitem{rodrigueswiki}
Wikipedia contributors, \emph{Rodrigues' rotation formula}, user-specified
statement, accessed November 14, 2026.
\url{https://en.wikipedia.org/wiki/Rodrigues%27_rotation_formula}.

\bibitem{proveit}
Vladimir Reshetnikov, \emph{ProveIt: ConciseIPRA}, repository snapshot
\texttt{c1bda76f7f86d2ade505355980826f625e36fff3}.
Inspected \texttt{README.md}, \texttt{PR.lean}, \texttt{Syntax.lean},
\texttt{BiInterpretation.lean}, \texttt{NativeBiInterpretation.lean},
and the root toolchain/build configuration.
\url{https://github.com/VladimirReshetnikov/ProveIt/tree/c1bda76f7f86d2ade505355980826f625e36fff3/Logic/PrimitiveRecursiveArithmetic/ConciseIPRA}.

\bibitem{drafts}
ProveIt research drafts, \texttt{docs/drafts/FTA\_iPRA.zip} and
\texttt{docs/drafts/triangle\_formulas\_ipra.zip}, at the same pinned
snapshot. Inspected their source discussions of PR-name schemas,
quantifier-free induction and formalization status. No geometric
result from these drafts is assumed here.
\url{https://github.com/VladimirReshetnikov/ProveIt/tree/c1bda76f7f86d2ade505355980826f625e36fff3/docs/drafts}.

\bibitem{bridgesrichman}
Douglas Bridges and Fred Richman, \emph{Varieties of Constructive
Mathematics}, London Mathematical Society Lecture Note Series 97,
Cambridge University Press, 1987, Chapter 1 (LLPO and constructive
existence). Publisher's excerpt:
\url{https://assets.cambridge.org/97805213/18020/excerpt/9780521318020_excerpt.pdf}.

\bibitem{kleene}
S.~C.~Kleene, ``On the interpretation of intuitionistic number theory,''
\emph{Journal of Symbolic Logic} \textbf{10}(4) (1945), 109--124.
\url{https://doi.org/10.2307/2269016}.
Author's article in an academic course archive:
\url{https://webspace.science.uu.nl/~ooste110/studsemIntMod/Kleene45.pdf}.

\bibitem{moschovakis}
Joan Rand Moschovakis and Garyfallia Vafeiadou, \emph{Intuitionistic
Mathematics and Logic}, authors' English translation of their 2009
chapter, section on number realizability and Nelson's theorem.
\url{https://nuprl-web.cs.cornell.edu/documents/Rand/intuitionistic-math2.pdf}.
\end{thebibliography}

\end{document}
